% id: 123-41
% book: 베이직쎈 미적분1 · 08 정적분
% section: 개념 51 구간에 따라 다르게 정의된 함수의 정적분
% passage: 다음 정적분의 값을 구하시오.
% figure: none
% answer: 
% answer_source: 
% variant_level: 0 (원본 전사)
% 이 파일은 build.mjs 가 items.json 에서 만든다. 고칠 때는 items.json 을 고친다.
\smash{\raisebox{6.5mm}{\dmkichul{베이직쎈 미적분1 123쪽 41번}}}
$f(x)=\begin{cases} 2x+1 & (x\ge 0) \\ x+1 & (x\le 0) \end{cases}$일 때, $\displaystyle\int_{-1}^{1} f(x)\mkern3mu dx$
\dmhint{$\displaystyle\int_{-1}^{1} f(x)\mkern3mu dx=\int_{-1}^{\text{\scriptsize\setlength{\fboxsep}{1.5pt}\framebox[1.5em]{\rule{0pt}{1.6ex}}}} f(x)\mkern3mu dx+\int_{\text{\scriptsize\setlength{\fboxsep}{1.5pt}\framebox[1.5em]{\rule{0pt}{1.6ex}}}}^{1} f(x)\mkern3mu dx$\\ $\displaystyle =\int_{-1}^{\text{\scriptsize\setlength{\fboxsep}{1.5pt}\framebox[1.5em]{\rule{0pt}{1.6ex}}}} (x+1)\mkern3mu dx+\int_{\text{\scriptsize\setlength{\fboxsep}{1.5pt}\framebox[1.5em]{\rule{0pt}{1.6ex}}}}^{1} (\blank[3.6em])\mkern3mu dx$\\ $\displaystyle =\Big[\dfrac{1}{2}x^2+x\Big]_{-1}^{\text{\scriptsize\setlength{\fboxsep}{1.5pt}\framebox[1.5em]{\rule{0pt}{1.6ex}}}}+\Big[\blank[3.6em]\Big]_{\text{\scriptsize\setlength{\fboxsep}{1.5pt}\framebox[1.5em]{\rule{0pt}{1.6ex}}}}^{1}$\\ $=\blank$}
